\subsection{A changing population of independent learners}
The implementation couples a shared physical world to private learner state
(Fig.~\ref{fig:architecture}). JSON specifies species, body settings, controller
architecture, actions, geography and feedback. Animals observe local state,
choose actions independently and receive individual feedback after a shared
transition. Resource growth, energy demand, attacks, starvation and birth change
the population through time. These are implemented capabilities, rather than
experimental evidence of emergent ecological competence.
\begin{figure}[t]
\centering
\begin{tikzpicture}[node distance=9mm and 18mm]
\node[block] (config) {Validated JSON\\species, maps, templates};
\node[block,right=of config] (world) {Python world\\bodies, resources, RNG};
\node[block,right=of world] (brains) {Private controllers\\parameters, buffers, RNGs};
\node[block,below=of world] (recorder) {Artifact writer\\replay + checkpoints\\no viewer required};
\node[block,below=of recorder] (files) {Saved artifacts\\replay, state, champions};
\node[block,left=of files] (server) {Read-only server\\requested files / chunks};
\node[block,left=of recorder] (viewer) {JavaScript viewer\\frames and metrics\\agent inspection};
\draw[arrow] (config)--(world);
\draw[arrow,transform canvas={yshift=2pt}] (world)--node[above,font=\scriptsize]{obs./reward}(brains);
\draw[arrow,transform canvas={yshift=-2pt}] (brains)--node[below,font=\scriptsize]{actions}(world);
\draw[arrow] (world)--(recorder);
\draw[arrow] (brains.south)|- (recorder.east);
\draw[arrow] (recorder)--(files);
\draw[arrow] (files)--(server);
\draw[arrow] (server)--(viewer);
\end{tikzpicture}
\caption{Runtime ownership. The training process writes artifacts; playback
reads them without a live world or learning controller. Arrows group data access,
not a new ordering of simultaneous actions.}\label{fig:architecture}
\end{figure}


\subsection{Spatial constraints and explicit sensory limits}
Water, elevation and obstacles constrain movement and energetic demand. Water
entry requires a swim decision; crossing is learned through the selected
controller rather than a separate route planner. Visibility uses distance and
wall occlusion, with an optional directional extension across water. Visible
targets are reduced to configurable top-$k$ slots per category, rather than all
agents being supplied as an unrestricted state vector
(Fig.~\ref{fig:environment-observation}). This separates physical accessibility,
visibility and the information actually exposed to a controller.
\begin{figure}[p]
\centering
\resizebox{\linewidth}{!}{%
\begin{tikzpicture}[x=6mm,y=6mm,font=\footnotesize]
\node[anchor=west,font=\bfseries] at (-.5,8.5) {a\quad Geography and visibility};
\fill[green!5] (-.5,-.5) rectangle (9.5,7.5);
\fill[cyan!20] (1.5,-.5) rectangle (2.5,7.5);
\draw[step=1,black!15,shift={(-.5,-.5)}] (0,0) grid (10,8);
\foreach \y in {0,1,2,4,5,6,7}{\fill[black!65] (4.5,\y-.5) rectangle (5.5,\y+.5);}
\draw[blue!75!black,dashed,thick] (1,3) circle (2);
\fill[blue!75!black] (1,3) circle (.19);
\node[anchor=east] at (.65,3) {$i$};
\fill[green!50!black] (4,3) circle (.18);
\draw[green!40!black,thick] (1,3)--(4,3);
\draw[red!75!black,thick,densely dashed] (1,3)--(7,5);
\fill[red!75!black] (7,5) circle (.18);
\node[anchor=west] at (7.3,5) {hidden};
\node[anchor=west] at (5.7,2.9) {gate};
\draw[-{Latex},cyan!50!black,thick] (2,6.6)--(2,5.6);
\draw[-{Latex},cyan!50!black,thick] (2,1.4)--(2,.4);
\node[anchor=north west,align=left,font=\scriptsize,text width=59mm] at (-.5,-1.1) {Dashed: base $R=2$; water bonus $B=1$.\\Green: $d=3$; one water cell extends sight.\\Red: walls still block extended sight.\\Water entry requires swim; arrows: current.};
\node[anchor=west,font=\bfseries] at (12,8.5) {b\quad Controller input: $v=3$, $k=4$};
\foreach \lo/\hi/\y/\txt/\col in {0/7/6.7/{Self and local movement}/blue,8/23/5.2/{Four resource slots}/green,24/39/3.7/{Four resource-eater slots}/teal,40/55/2.2/{Four prey-eater slots}/orange,56/67/.7/{Local terrain}/cyan,68/77/-0.8/{Mating readiness and nearest mate}/violet}{
 \filldraw[fill=\col!10,draw=\col!55!black] (12,\y-.5) rectangle (26,\y+.5);
 \node[anchor=west,font=\scriptsize] at (12.2,\y) {\lo--\hi\quad\txt};
}
\draw[-{Latex},black!65] (10,3.7)--(11.7,3.7);
\node[font=\scriptsize,align=center] at (11,5) {filter\\rank\\pad};
\node[anchor=west] at (12,-2.3) {Each target slot: $(\Delta x,\Delta y,d,\mathrm{presence})$.};
\node[anchor=west] at (12,-3.1) {Missing slots are zero; total width is 78.};
\end{tikzpicture}}
\caption{\textbf{Geography, sensing and the packed observation vector.}
\textbf{a}, Controlled schematic, not an experimental replay. The water bonus
can expose a target beyond the base circle but does not remove wall occlusion.
The drawing illustrates sensing; movement eligibility remains a separate check.
\textbf{b}, Zero-based input ranges for observation version 3 with four targets
per category. Each category is independently filtered and ranked. Terrain and
mating fields are separate from general target slots; no target brain or health
is disclosed. The layout does not imply a convolutional image input.}
\label{fig:environment-observation}
\end{figure}


\subsection{Reproductive inheritance and independent descendants}
Birth requires compatible nearby animals to choose reproduction in the same
step and satisfy physical eligibility. Both parents contribute parameters;
rare mutation affects only the new child. The child receives fresh optimizer,
experience and random-generator state (Fig.~\ref{fig:inheritance}). Between-run
champion selection is a separate mechanism: it retains strictly improving
species/sex-specific models according to the selected criterion, including
lifespan, feeding, kills or founder-family lineage. Such rankings summarize
training history and require held-out tests before being interpreted as fitness.
\begin{figure}[p]
\centering
\begin{tikzpicture}[x=1cm,y=1cm,font=\footnotesize]
\node[anchor=west,font=\bfseries] at (0,6.8) {a\quad Two-parent parameter inheritance};
\node[anchor=east] at (1.6,5.8) {Parent $p$};
\node[anchor=east] at (1.6,4.9) {Parent $q$};
\foreach \i in {0,...,7}{
 \filldraw[fill=blue!25,draw=white] (1.8+\i*.45,5.55) rectangle +( .45,.5);
 \filldraw[fill=orange!35,draw=white] (1.8+\i*.45,4.65) rectangle +( .45,.5);
}
\draw[arrow] (3.6,4.5)--(3.6,3.9);
\node[anchor=east] at (1.6,3.5) {Child};
\foreach \i/\col in {0/blue,1/orange,2/orange,3/blue,4/blue,5/orange,6/blue,7/orange}{
 \filldraw[fill=\col!30,draw=white] (1.8+\i*.45,3.25) rectangle +(.45,.5);
}
\draw[red!70!black,thick] (4.5,3.25) rectangle +(.45,.5);
\node[anchor=west,font=\scriptsize,align=left] at (5.7,3.5) {Rare child-only\\weight perturbation};
\node[draw=black!50,rounded corners,fill=black!3,align=center,text width=6.2cm,minimum height=12mm] (fresh) at (3.7,1.8) {Fresh optimizer, empty experience,\\independent random streams and temporal state};
\draw[arrow] (3.6,3.1)--(fresh.north);
\node[align=left,anchor=west] at (.2,.35) {Paired decisions and physical eligibility are required.\\Only parameters are mixed; training state is private.};
\node[anchor=west,font=\bfseries] at (9,6.8) {b\quad Founder-family membership};
\foreach \id/\x/\y/\txt in {A/9.7/5.7/{A},B/13.7/5.7/{B},C/10.5/4.2/{C},D/12.9/4.2/{D},E/11.7/2.7/{E}}{
 \node[circle,draw=black!60,fill=green!8,minimum size=7mm] (\id) at (\x,\y) {\txt};
}
\foreach \from/\to in {A/C,A/D,B/C,B/D,C/E,D/E}{\draw[arrow] (\from)--(\to);}
\node[anchor=west,align=left] at (9,1.3) {$C,D,E$ have founder roots $\{A,B\}$.\\Family A: $A,C,D,E$; family B: $B,C,D,E$.};
\node[anchor=west,align=left] at (9,.1) {An $E$ contribution counts once per family.\\Rank uses the strongest family, not their sum.};
\end{tikzpicture}
\caption{\textbf{Parameter inheritance and ancestry are different structures.}
\textbf{a}, Colors indicate parental parameter origin, not reproductive sex;
cell contents are schematic. Uniform crossover follows parental learning for
that step. Mutation can perturb selected child matrix weights, while training
state is freshly initialized. \textbf{b}, A shared descendant belongs to both
founder families without duplicate membership within either family. The diagram
illustrates ancestry bookkeeping, not measured reproductive success. Personal
eligibility and the complete ranking tuple are described in Methods.}
\label{fig:inheritance}
\end{figure}


\subsection{Artifacts preserve distinct kinds of state}
Replay, exact continuation and model transfer require different data
(Fig.~\ref{fig:artifacts}). A viewer accesses independently decodable replay
chunks without restoring private controllers. A checkpoint restores complete
learning state. Champion metadata references compatible model blobs shared by
content hash, permitting parameter transfer without duplicating unchanged
weights. These boundaries allow headless simulation and later inspection while
limiting the active replay working set.
\begin{figure}[t]
\centering
\begin{tikzpicture}[node distance=7mm and 8mm]
\node[block] (replay) {Replay chunk\\keyframe + field deltas};
\node[block,right=of replay] (checkpoint) {Exact checkpoint\\metadata + brain records};
\node[block,right=of checkpoint] (champion) {Champion metadata\\rank + source + references};
\node[block,below=of replay] (view) {Viewer access\\one world chunk at a time};
\node[block,below=of checkpoint] (restore) {Restore all learning state\\models, optimizer, RNG, buffers};
\node[block,below=of champion] (pool) {Shared model blob\\signature + precision + arrays};
\node[block,below=of pool] (fresh) {Fresh controller initialization\\copy compatible parameters};
\draw[arrow] (replay)--(view);
\draw[arrow] (checkpoint)--(restore);
\draw[arrow] (champion)--(pool);
\draw[arrow] (pool)--(fresh);
\end{tikzpicture}
\caption{Artifact structure determines access semantics. Lossless compression
applies to each binary payload, while readable metadata remains outside it.
Checkpoint loading materializes full learning state; champion loading and
viewer replay serve different purposes.}\label{fig:artifacts}
\end{figure}

